\subsubsection{Momentos de inércia}
\label{sec:dinAcoplada_momentos_inercia}

A matriz de inércia generalizada $M(\vec q\,)$ é introduzida a partir da
forma quadrática da energia cinética, conforme
\eqref{eq:T_matriz_corrigida} e \eqref{eq:M_total_formula_corrigida} e em
linha com \cite{spong2006robot,featherstone2008}:

\begin{equation}
T =
\frac{1}{2}\,\dot{\vec q}^{\,T}\,M(\vec q\,)\,\dot{\vec q}.
\end{equation}

Os elementos $m_{kj}(\vec q\,)$ de $M(\vec q\,)$ têm interpretação direta:
os termos da diagonal representam as inércias equivalentes associadas a
cada grau de liberdade, enquanto os termos fora da diagonal expressam os
acoplamentos inerciais entre a atitude da base e as juntas do
manipulador \cite{Geradin2004}.

Adota-se o vetor de coordenadas generalizadas conforme
\eqref{eq:q_generalizado_base_mr},

\begin{equation}
\vec q
=
\begin{bmatrix}
\vec\eta_B \\
\vec\theta
\end{bmatrix},
\end{equation}

onde $\vec\eta_B$ descreve a atitude da base e $\vec\theta$ reúne as
coordenadas articulares do \gls{mr}. Com essa escolha, a matriz
$M(\vec q\,)$ assume a forma particionada usual para sistemas de base
flutuante \cite{featherstone2008,Wilde2018}:

\begin{equation}
M(\vec q\,) =
\begin{bmatrix}
M_{BB}(\vec q\,) & M_{BM}(\vec q\,) \\
M_{MB}(\vec q\,) & M_{MM}(\vec q\,)
\end{bmatrix},
\label{eq:M_blocos}
\end{equation}

em que $M_{BB}$ agrupa as inércias equivalentes associadas à atitude da
base física, $M_{MM}$ contém as contribuições inerciais do manipulador e
os blocos $M_{BM}$, $M_{MB}$ representam os termos de acoplamento
inercial entre base e \gls{mr}.

Do ponto de vista construtivo, a expressão
\eqref{eq:M_total_formula_corrigida} mostra que $M(\vec q\,)$ é obtida
somando-se as contribuições espaciais de cada elo e da base, isto é,

\begin{equation}
M(\vec q\,) =
\sum_{i=1}^{n}
J_{\text{sp},i}^{T}(\vec q\,)\,
M_{i,\text{spatial}}\,
J_{\text{sp},i}(\vec q\,)
+
M_{\text{base},\text{spatial}},
\end{equation}

onde $J_{\text{sp},i}(\vec q\,)$ é o jacobiano espacial do elo $i$ e
$M_{i,\text{spatial}}$ é o operador de inércia espacial associado à
massa e ao tensor de inércia do elo no formalismo 6D
\cite{featherstone2008}. Essa construção garante a consistência entre a
cinemática diferencial (via jacobianos espaciais) e a estrutura da
matriz de inércia generalizada empregada nas equações de Lagrange em
\eqref{eq:dinamica_geral}.
